\chapter{La repeated sprint ability nel calcio}

% ---------------------------------------------------------------------------
\section{Che cos'è la repeated sprint ability}

\textbf{Repeated sprint ability} (RSA): la capacità, nell'ambito di alcune discipline sportive come
il calcio, di effettuare \textbf{sprint massimali} alternati a \textbf{periodi di recupero}, che
possono essere completi o consistere in attività di bassa intensità.

\rubrica{Quanto dura uno sprint}
\begin{itemize}
  \item molti autori considerano RSA anche sprint della durata di \textbf{30 secondi e oltre};
  \item in ambito calcistico è più adeguato limitarsi agli articoli che individuano la RSA in sprint
    la cui \textbf{durata massima} sia dell'ordine di \textbf{circa 10 secondi e non oltre}.
\end{itemize}

\subsection{\texorpdfstring{\textit{Intermittent} e \textit{repeated sprint exercise}}{Intermittent e repeated sprint exercise}}

Due sigle da non confondere. Ciò che le distingue in modo netto è la \textbf{durata del recupero};
in entrambi i casi lo sprint è di breve durata.

\begin{itemize}
  \item \textbf{ISE} (\textit{intermittent sprint exercise}): sprint di corta durata separati da
    \textbf{recuperi completi}, superiori ai \textbf{60 secondi};
  \item \textbf{RSE} (\textit{repeated sprint exercise}): sprint di breve durata separati da
    \textbf{brevi recuperi}, generalmente inferiori ai \textbf{60 secondi}.
\end{itemize}

La conseguenza pratica è netta:

\begin{itemize}
  \item nell'\textbf{ISE} si registra un decremento della performance \textbf{limitato}, e a volte
    nemmeno quello;
  \item nell'\textbf{RSE} si assiste a un decremento della prestazione \textbf{marcato}.
\end{itemize}

% ---------------------------------------------------------------------------
\section{I fattori che limitano la prestazione}

\subsection{La fatica muscolare}

La \textbf{fatica muscolare} --- o \textbf{periferica}, con termine più tecnico --- è un fattore in
grado di compromettere la performance durante una serie di RSE. Molti ricercatori hanno dimostrato
che nel corso di una serie di sprint ripetuti si verificano \textbf{modifiche nella conduzione dello
stimolo nervoso}, che rifletterebbero cambiamenti a carico di particolari ioni --- \textbf{sodio} e
\textbf{potassio} --- che regolano l'attività muscolare.

\subsection{La fosfocreatina}

Il primo intervento è inserire nel programma esercitazioni rivolte all'incremento della
\textbf{resistenza alla forza esplosiva}; \textbf{non basta}, però, perché la RSA dipende da molti
fattori interdipendenti e non si riduce a uno scadimento delle caratteristiche di forza.

La \textbf{deplezione di fosfocreatina} (PCr) può costituire un fattore \textbf{predominante} nella
diminuzione della performance durante RSE con recupero \textbf{minore o uguale a 30 secondi}
(Bangsbo e coll., 1992).

\subsection{Il glicogeno muscolare}

\begin{itemize}
  \item numerosi studi hanno dimostrato la relazione fra la capacità di mantenere un'alta produzione
    di potenza muscolare durante sforzi di \textbf{breve durata} (minori di 30 secondi), ripetuti in
    modalità intermittente, e la \textbf{disponibilità di glicogeno muscolare} (Casey e coll., 1996;
    Balsom e coll., 1999);
  \item sul contributo del meccanismo anaerobico alla prestazione di sprint --- e quindi sulla sua
    capacità di produrre ATP --- c'è \textbf{sostanziale unanimità} nel considerare il meccanismo
    \textbf{glicolitico} un fattore sostanziale nell'ambito della RSA.
\end{itemize}

\subsection{I sistemi di ripristino energetico}

Il peso dei tre meccanismi cambia con la durata dell'accelerazione o dello sprint:

\begin{table}[htbp]
  \centering
  \small
  \renewcommand{\arraystretch}{1.3}
  \begin{tabularx}{\textwidth}{|X|c|c|c|}
    \hline
    \textbf{Durata di accelerazioni e/o sprint} & \textbf{Aerobico} &
      \textbf{Anaerobico lattacido} & \textbf{Anaerobico alattacido} \\
    \hline
    30 secondi & 38\% & 45\% & 17\% \\
    \hline
    Tra 22 e 11 secondi & 30\% & 47\% & 23\% \\
    \hline
    Minore di 10 secondi & circa 10\% & 45\% & 50\% \\
    \hline
    Tra 3 e 1,5 secondi & non significativo & tra 50\% e 20\% & tra 50\% e 80\% \\
    \hline
  \end{tabularx}
\end{table}

Al ridursi della durata il contributo \textbf{alattacido} cresce fino a diventare dominante e quello
\textbf{aerobico} si annulla, mentre il \textbf{lattacido} resta rilevante su tutto l'arco.

\subsection{La disponibilità di ossigeno}

\begin{itemize}
  \item i soggetti con un elevato $VO_2max$ e/o con una \textbf{cinetica del consumo di ossigeno}
    particolarmente efficiente mostrano una maggiore \textbf{resistenza alla fatica} (Bishop e
    coll., 2004; Dupont e coll., 2005);
  \item l'intervento del meccanismo aerobico \textbf{aumenta parallelamente alla durata} dello
    sprint, e diminuisce al ridursi di questa;
  \item in un RSE l'intervento del meccanismo aerobico aumenterebbe già del \textbf{15\% al secondo
    sprint}, contenendo la perdita di performance sull'intera serie.
\end{itemize}

\subsection{L'acidosi, e perché è contestata}

L'\textbf{accumulo di lattato} è stato da sempre indicato come uno dei fattori maggiormente in grado
di interferire con i meccanismi di contrazione muscolare, e quindi come determinante nella
prestazione basata sulla RSA (Bishop e coll., 2004; Allen e coll., 2008).

Queste ipotesi sono però oggetto di contestazione per almeno tre motivi:

\begin{enumerate}
  \item il recupero della \textbf{forza} e della \textbf{potenza} muscolare sarebbe \textbf{più
    rapido} del ritorno a valori di pH normali --- cioè di concentrazione di ioni idrogeno, e quindi
    di acidosi: in altre parole il recupero di forza e potenza sarebbe abbastanza
    \textbf{indipendente} dallo smaltimento del lattato prodotto;
  \item in alcuni lavori sperimentali si è registrata un'importante produzione di potenza muscolare
    \textbf{anche in presenza di forte acidosi} dei muscoli coinvolti;
  \item l'ingestione di \textbf{bicarbonato di sodio} --- sostanza nota per il suo potere tampone,
    cioè per la capacità di inattivare gli ioni idrogeno responsabili dell'acidosi --- \textbf{non ha
    alcun effetto positivo} sulla RSA (Glaister, 2005).
\end{enumerate}

\subsection{Il quadro d'insieme}

\begin{table}[htbp]
  \centering
  \small
  \renewcommand{\arraystretch}{1.3}
  \begin{tabularx}{\textwidth}{|X|X|}
    \hline
    \textbf{Fattore} & \textbf{Influenza sulla RSA} \\
    \hline
    Conduzione dello stimolo nervoso & In grado di influenzare la performance \\
    \hline
    Disponibilità di PCr & Dubbio \\
    \hline
    Glicogeno muscolare & In grado di influenzare la performance \\
    \hline
    Livello di performance aerobica & In grado di influenzare la performance \\
    \hline
    Acidosi & Per alcuni aspetti contestabile \\
    \hline
  \end{tabularx}
\end{table}

% ---------------------------------------------------------------------------
\section{L'allenamento della RSA}

\subsection{I due mezzi principali}

\rubrica{RSA breve}
La variante più semplice:

\begin{itemize}
  \item \textbf{6--10 ripetizioni} a intensità massimale, a navetta o con altri tipi di cambio di
    direzione;
  \item distanze da \textbf{10 $+$ 10\,m} a \textbf{20 $+$ 20\,m};
  \item recupero di \textbf{20--40$''$};
  \item il tutto organizzato in \textbf{2--3 serie}.
\end{itemize}

Il rapporto fra \textbf{fase attiva} e \textbf{pausa} governa il tipo di sollecitazione: minore è il
rapporto --- p.\,es.\ $5''$ di attività con $30''$ di pausa danno un rapporto di \textbf{1:6} ---
\textbf{maggiore} è la sollecitazione \textbf{muscolare} e minore quella \textbf{metabolica};
viceversa al crescere del rapporto.

\rubrica{RSA based training}
Variante dell'RSA breve:

\begin{itemize}
  \item azioni alla massima intensità di circa \textbf{6--7$''$} (p.\,es.\ navetta 20 $+$ 20\,m);
  \item rapporto con pausa \textbf{da fermo} di \textbf{1:3}, cioè $20''$;
  \item \textbf{6 ripetizioni} per serie, per un totale di \textbf{240--300\,m} a serie;
  \item \textbf{3 serie} complessive.
\end{itemize}

\rubrica{Il confronto con l'interval training (Ferrari e coll., 2008)}
L'RSA based training è stato confrontato con le ripetute classiche, $4 \times 4'$ al
\textbf{90--95\% della $FC_{max}$}. Dopo \textbf{7 settimane}, due volte a settimana:

\begin{itemize}
  \item le due metodologie offrivano gli \textbf{stessi miglioramenti} di $VO_2max$;
  \item l'RSA based training ha però permesso \textbf{maggiori incrementi} nella distanza dello
    \textbf{Yo-Yo Intermittent Recovery Test}.
\end{itemize}

\subsection{Dove collocarlo nel ciclo di allenamento}

\begin{itemize}
  \item \textbf{prerequisito}: i giocatori devono avere già appreso correttamente i \textbf{cambi di
    direzione}, cioè averci lavorato in precedenza in modo analitico, altrimenti non sono in grado di
    eseguirli alla massima intensità;
  \item la richiesta di attività massimali esige che il lavoro venga somministrato solo in
    condizioni di \textbf{freschezza muscolare};
  \item è \textbf{sconsigliabile} effettuarlo a ridosso delle partite;
  \item di conseguenza va collocato nella \textbf{parte centrale della settimana}, in giorni in cui
    non si lavora sulla potenza muscolare.
\end{itemize}

\subsection{Le esercitazioni di resistenza speciale alla forza}

\begin{itemize}
  \item esercitazioni a \textbf{intensità massimale} con sollecitazioni di tipo anaerobico, con cui
    migliorare la \textbf{capacità di recupero in tempi brevi}: scatti di lunghezza diversa,
    intercalati da cambi di direzione con \textbf{angolazioni diverse};
  \item la muscolatura viene impegnata \textbf{eccentricamente} nelle fasi di decelerazione e
    \textbf{concentricamente} nelle accelerazioni, con reiterati impegni di \textbf{forza
    esplosiva}.
\end{itemize}

\rubrica{Distanze e volumi}
\begin{itemize}
  \item \textbf{20--40\,m} nella singola ripetizione;
  \item quantità per seduta fra i \textbf{300 e gli 800\,m}, in funzione dello \textit{status} fisico
    degli atleti.
\end{itemize}

\rubrica{La progressione del carico}
L'evoluzione deve rispettare il criterio di \textbf{progressività}:

\begin{enumerate}
  \item prima si agisce sulla \textbf{quantità}: aumento del numero di ripetizioni e del numero di
    cambi di direzione;
  \item poi sulla \textbf{densità}: diminuzione dei tempi di recupero.
\end{enumerate}

Resta fermo che l'\textbf{intensità} deve attestarsi sempre sui massimi livelli esprimibili dai
calciatori in quel particolare momento.

\subsection{Sprint training e decelerazioni}

\begin{itemize}
  \item diversi autori sostengono che la performance in prove di \textbf{sprint lineari} e in quelle
    \textbf{con cambi di direzione} richieda capacità specifiche \textbf{differenti}: è opportuno
    quindi inserire nel programma anche accelerazioni lineari su distanze fra \textbf{5 e 30\,m};
  \item dopo l'accelerazione massimale si richiede all'atleta una \textbf{decelerazione repentina},
    che arresti la corsa nel minor tempo possibile, sollecitando gli arti inferiori a elevati
    \textbf{gradienti di forza in regime eccentrico};
  \item così si indirizza la capacità di \textbf{accelerazione-decelerazione} verso un regime di
    \textbf{resistenza}, stimolando i meccanismi attivi durante le pause e incrementando la capacità
    di recupero dopo lo sforzo massimale.
\end{itemize}

\subsection{Le accelerazioni in salita}

\begin{itemize}
  \item già nel \textbf{1984} Arcelli affermava che le accelerazioni condotte in salita permettono
    adattamenti utili allo sviluppo della condizione fisica del calciatore;
  \item la maggior parte delle pubblicazioni internazionali usa protocolli con sforzi dell'ordine di
    una \textbf{decina di secondi e più};
  \item comunemente si utilizzano però sprint su distanze inferiori: \textbf{10--30\,m},
    \textbf{3--6 secondi}, con una \textbf{pendenza fra il 7 e il 10\%}.
\end{itemize}

\rubrica{Come tarare la salita (Vittori)}
Testare prima il \textit{best trial} sulla \textbf{distanza in piano} --- p.\,es.\ 30\,m in $4''$ ---
e scegliere poi una salita che provochi un peggioramento di \textbf{80 centesimi}, cioè $4''80$.

\rubrica{Che cosa cambia rispetto al piano}
\begin{itemize}
  \item il \textbf{pattern di attivazione muscolare} è significativamente differente: sono
    sollecitati in misura \textbf{superiore} glutei, gruppo dei vasti e soleo, in misura
    \textbf{inferiore} gracile, semitendinoso e semimembranoso;
  \item dal punto di vista \textbf{cinematico} l'elemento principale è la \textbf{riduzione
    dell'ampiezza del passo}: una particolarità che può risultare favorevole per ridurre il rischio
    di infortunio ai \textbf{muscoli flessori} della gamba.
\end{itemize}

\subsection{L'allenamento della forza}

\begin{itemize}
  \item le varie espressioni di forza sono ritenute fra i \textbf{fattori principali} che influenzano
    le performance fisiche;
  \item nel calcio le espressioni fondamentali sono quelle che richiedono una \textbf{velocità di
    accorciamento elevata} e sono modulate dal sistema neuromuscolare: la \textbf{forza esplosiva} e
    la \textbf{forza dinamica massima};
  \item la programmazione di allenamenti mirati a queste due espressioni deve quindi avere notevole
    rilevanza ai fini prestativi calcistici.
\end{itemize}

% ---------------------------------------------------------------------------
\section{Che cosa dice la ricerca}

\subsection{Il problema dei protocolli (Spencer e coll.)}

I protocolli usati dai diversi autori \textbf{non corrispondono ai modelli prestativi} degli sport di
squadra: la durata degli sforzi è troppo lunga e i recuperi troppo brevi.

\begin{table}[htbp]
  \centering
  \small
  \renewcommand{\arraystretch}{1.3}
  \begin{tabularx}{\textwidth}{|X|c|c|}
    \hline
     & \textbf{Sprint} & \textbf{Recuperi} \\
    \hline
    Negli \textbf{sport di squadra} & $< 3$\,s & $> 60$\,s \\
    \hline
    Nei \textbf{test} proposti dalle ricerche & $> 6$\,s & $< 30$\,s \\
    \hline
  \end{tabularx}
\end{table}

Nella sua \textit{review} Spencer rileva inoltre che molte ricerche sono state realizzate con un
\textbf{cicloergometro}, che non riproduce le modalità esecutive della \textbf{corsa in
accelerazione}.

\subsection{Giocatori veloci e giocatori lenti (Méry e Cometti, 2004)}

Ricerca dell'Università di \textbf{Digione} sulla diversità di risposta alle serie di sprint ripetuti
fra giocatori \textbf{veloci} e giocatori più \textbf{lenti}. Le serie di riferimento sono quelle
eseguite con pause di \textbf{20 secondi}, e i risultati sono letti su blocchi di cinque sprint
($S_{1\text{-}5}$, $S_{6\text{-}10}$, $S_{11\text{-}15}$, $S_{16\text{-}20}$).

\rubrica{I tempi sui 15 metri}
\begin{itemize}
  \item i giocatori \textbf{lenti} mostrano un \textbf{aumento significativo} dei tempi --- che
    peggiorano rapidamente e poi si stabilizzano ($p < 0{,}05$ dal secondo blocco in poi): non
    riescono a mantenere la performance iniziale;
  \item i giocatori \textbf{veloci} subiscono un \textbf{lieve incremento non significativo}: sono in
    grado di mantenere la performance.
\end{itemize}

\rubrica{L'andamento della velocità}
\begin{itemize}
  \item nei \textbf{lenti} la velocità misurata a ciascun appoggio \textbf{diminuisce nettamente} nel
    corso delle ripetizioni ($p < 0{,}05$ al secondo blocco, $p < 0{,}01$ al terzo e al quarto);
  \item nei \textbf{veloci} non si osservano differenze significative fra i primi sprint e i
    successivi: non perdono velocità, e dimostrano quindi di essere più \textbf{resistenti alla
    velocità}.
\end{itemize}

\rubrica{Il passo successivo: che cosa allenare nei lenti}
La conclusione parrebbe che i soggetti lenti debbano allenare la \textbf{velocità} e non la
resistenza alla velocità. Ma la velocità dipende da due fattori --- \textbf{frequenza} e
\textbf{lunghezza del passo} --- e i risultati mostrano un calo della \textbf{frequenza}. Di qui una
seconda indagine:

\begin{itemize}
  \item \textbf{3 settimane} di allenamento, \textbf{3 sessioni} a settimana;
  \item un gruppo (\textit{sprint-frequenza}) con esercizi specifici per migliorare la
    \textbf{frequenza dei passi};
  \item l'altro con un lavoro per aumentare la \textbf{lunghezza dei passi}, attraverso un
    allenamento di \textbf{forza massimale}.
\end{itemize}

\rubrica{Il risultato}
Confrontando le curve prima e dopo le tre settimane nel gruppo ``potenziamento'', la performance
peggiora in modo significativo a partire dalla \textbf{3\textsuperscript{a} ripetizione} nei test
effettuati \textbf{prima}, e solo dalla \textbf{14\textsuperscript{a}} in quelli effettuati
\textbf{dopo}: l'allenamento della forza massimale sposta molto più in là il punto in cui la fatica
comincia a incidere.

\rubrica{La conclusione}
Un allenamento di potenziamento muscolare orientato all'incremento della \textbf{forza massimale}
sembra in grado di migliorare la capacità di \textbf{resistenza alla velocità} in modo più efficace
rispetto a uno orientato al miglioramento della \textbf{frequenza dei passi}.

% ---------------------------------------------------------------------------
\section{Le esercitazioni con la palla ad elevata intensità}

Le esercitazioni con la palla ad alta intensità vengono sempre più utilizzate per integrare gli
elementi della preparazione fisica della squadra, e quando sono proposte \textbf{rispettando i
criteri corretti} sono realmente efficaci a questo scopo (Impellizzeri e coll., 2006). Le risposte
fisiologiche variano però in funzione di tre fattori.

\subsection{Il numero dei giocatori}

\begin{itemize}
  \item i giocatori realizzano gesti tecnici con una \textbf{frequenza correlata al tipo di
    esercitazione}: p.\,es.\ nell'$1 < 1$ il numero di \textit{tackle} è superiore rispetto alle
    altre mini-partite;
  \item determinante è il numero di \textbf{possessi palla} per ciascun giocatore: Balsom (1999) ha
    osservato che nel $3 < 3$ è \textbf{tre volte superiore} rispetto al $7 < 7$.
\end{itemize}

\subsection{La dimensione dello spazio di gioco}

Rampinini e coll.\ (2007) hanno confermato che la dimensione dello spazio ha effetto sulla
\textbf{frequenza cardiaca}, sulla \textbf{concentrazione ematica del lattato} e sul \textbf{livello
di fatica percepito}, in esercitazioni $3 < 3$, $4 < 4$, $5 < 5$ e $6 < 6$.

\subsection{L'incitamento da parte dell'allenatore}

È l'elemento che più spesso non viene considerato a sufficienza, e che influenza notevolmente
l'intensità dell'esercizio e quindi gli effetti allenanti:

\begin{itemize}
  \item quando il tecnico \textbf{incita costantemente} i propri atleti, l'intensità media
    dell'esercitazione risulta \textbf{significativamente superiore};
  \item la \textbf{frequenza cardiaca} è correlata in maniera \textbf{lineare} con l'intensità
    dell'esercizio, mentre la concentrazione del \textbf{lattato} varia in modo
    \textbf{esponenziale}.
\end{itemize}

% ---------------------------------------------------------------------------
\section{La valutazione della RSA}

\subsection{Il test navetta di Capanna}

È il test a intensità massimale più conosciuto per indagare le capacità lattacide dei calciatori:
\textbf{6 navette di 20\,m} ($20 + 20$) con $20''$ di recupero, elaborate come valore medio delle sei
prove o come decremento percentuale fra le prime due e le ultime due (Rampinini, 2006). Protocollo,
elaborazione e parametri di riferimento per i professionisti sono già stati esposti fra i test di
valutazione funzionale (§\ref{sec:test-capanna}).

\subsection{L'High Intensity Test (HIT)}

Protocollo ideato dallo \textbf{Human Performance Lab Mapei} per misurare alcune risposte
fisiologiche all'esercizio condotto ad alta intensità, basato su corse a navetta da eseguire a
velocità \textbf{submassimale}.

\rubrica{Protocollo}
\begin{itemize}
  \item $10 \times 10''$ di corsa a \textbf{18\,km/h} su una distanza di \textbf{25 $+$ 25\,m};
  \item \textbf{recupero passivo};
  \item si misurano nel sangue la concentrazione di \textbf{lattato} (La), gli \textbf{ioni
    idrogeno} e il grado di acidità (\textbf{pH}).
\end{itemize}

\rubrica{Che cosa valuta}
\begin{itemize}
  \item la funzionalità \textbf{congiunta} dei sistemi tampone del muscolo e del sangue e del
    meccanismo aerobico, che vengono stimolati in maniera \textbf{integrata} per fronteggiare le
    richieste di questo tipo di sforzo;
  \item le risposte fisiologiche all'esercizio intermittente di intensità submassimale sembrano
    rappresentare un insieme di \textbf{indicatori utili} per la misura delle capacità prestative del
    giocatore di calcio;
  \item per come è costruito è però più un test da \textbf{laboratorio} che da campo.
\end{itemize}
